\documentclass[12pt,a4paper]{article}
\usepackage[utf8]{inputenc}
\usepackage[T1]{fontenc}
\usepackage[polish]{polski}
\usepackage{amsmath}
\usepackage{amsfonts}
\usepackage{amssymb}
\usepackage{booktabs}
\usepackage{geometry}
\usepackage{enumitem}

\geometry{margin=2cm}

\title{Systemy czasu rzeczywistego, Projekt 5 -- Baza teoretyczna}
\author{Piotr Gugnowski, wydział EiTI, nr indeksu 320690}
\date{Czerwiec 2026}

\pagestyle{empty}

\begin{document}

\maketitle

% ===============================================================================
\section*{A.1: Zmiennoprzecinkowa reprezentacja liczb binarnych}
% ===============================================================================

Każdą niezerową liczbę rzeczywistą można zapisać w~\textbf{postaci znormalizowanej}:
\[
  x = (-1)^S \times 1{,}M \times 2^E
\]
gdzie:
\begin{itemize}
  \item $S \in \{0, 1\}$: bit znaku ($0$ = dodatnia, $1$ = ujemna),
  \item $1{,}M$: mantysa znormalizowana; część całkowita wynosi zawsze~1
        (tzw.\ \textit{implicit leading bit} lub \textit{hidden bit}),
  \item $E$: wykładnik całkowity (może być ujemny).
\end{itemize}

Mantysa $M$ to część ułamkowa: $M = \sum_{k=1}^{p} b_k \cdot 2^{-k}$,
gdzie $p$ to liczba bitów precyzji.

\medskip
\noindent\textbf{Intuicja.}
Zapis znormalizowany to po prostu przesunięcie przecinka binarnego tak,
aby przed nim stała dokładnie jedna cyfra \texttt{1}.
Wykładnik $E$ mówi o~ile pozycji i~w~którą stronę przesunięto przecinek
(ujemny $E$ -- w~lewo, dodatni -- w~prawo).
Mantysa $M$ to bity stojące \emph{za} tym przecinkiem -- cyfry po części całkowitej.

Przykład krok po kroku dla liczby~$9$:
\begin{enumerate}
  \item Zapisujemy $9$ binarnie: $9 = 1001_2$.
  \item Przesuwamy przecinek o~3~miejsca w~lewo, żeby zostało $1{,}\ldots$:
        $9 = 1{,}001_2 \times 2^3$.
  \item Stąd $E = 3$, a $M = 001\underbrace{00\cdots0}_{20}$ (23~bity łącznie).
  \item Sprawdzenie: $1{,}001_2 = 1 + 0{\cdot}2^{-1} + 0{\cdot}2^{-2} + 1{\cdot}2^{-3}
        = 1 + \tfrac{1}{8} = 1{,}125$;\quad
        $1{,}125 \times 2^3 = 1{,}125 \times 8 = 9$. $\checkmark$
\end{enumerate}

% ===============================================================================
\section*{A.2: Standard IEEE~754 single precision (32-bit)}
% ===============================================================================

\begin{center}
\begin{tabular}{|c|c|c|}
\hline
\textbf{Bit 31} & \textbf{Bity 30--23} & \textbf{Bity 22--0} \\
\hline
znak $S$ & wykładnik ze stałą $E_b$ (8 bitów) & mantysa $M$ (23 bity) \\
\hline
\end{tabular}
\end{center}

Wykładnik przechowywany jest z~\textbf{nadmiarem (biasem) 127}:
$E_b = E_{\text{actual}} + 127$, co pozwala zapisać ujemne wykładniki
bez bitu znaku. Zarezerwowane wartości $E_b$:

\begin{center}
\begin{tabular}{lll}
\toprule
$E_b$ & $M$ & Znaczenie \\
\midrule
$1\ldots254$ & dowolna & liczba znormalizowana: $(-1)^S \times (1{+}M) \times 2^{E_b - 127}$ \\
\bottomrule
\end{tabular}
\end{center}

Zakres liczb znormalizowanych dodatnich:
$\bigl[2^{-126},\ (2 - 2^{-23}) \cdot 2^{127}\bigr]
\approx \bigl[1{,}18 \times 10^{-38},\ 3{,}40 \times 10^{38}\bigr]$.

% ===============================================================================
\section*{A.3: Kod "wewnętrzny"}
% ===============================================================================

\begin{center}
\begin{tabular}{|c|c|c|c|}
\hline
\textbf{Bit 31} & \textbf{Bity 30--24} & \textbf{Bit 23} & \textbf{Bity 22--0} \\
\hline
"1" (jawny bit wiodący) & wykładnik $E$ bez stałej (7 bitów) & "0" (padding) & mantysa $M$ (23 bity) \\
\hline
\end{tabular}
\end{center}

Wartość liczby:
\[
  x = (1 + M \cdot 2^{-23}) \times 2^E
  \qquad E \in [0,\,127],\quad M \in [0,\,2^{23}-1]
\]

\medskip
\noindent\textbf{Co oznacza każde pole w~praktyce:}
\begin{itemize}
  \item \textbf{Bit~31 = 1 (jedynka wiodąca):} w~IEEE~754 bit wiodący mantysy
        ($1{,}M$) jest \emph{niejawny} -- procesor zakłada, że tam jest \texttt{1},
        ale jej nie zapisuje. Kod wewnętrzny zapisuje ją \emph{jawnie},
        więc bit~31 zawsze wynosi~\texttt{1} dla liczb niezerowych.
  \item \textbf{Bity 30--24 -- wykładnik $E$:} mówi przez jaką potęgę dwójki
        mnożymy mantysę. Dla $9 = 1{,}001_2 \times 2^3$ jest $E = 3$.
        W~IEEE~754 wykładnik jest przesunięty o~127, tutaj jest zapisany wprost.
  \item \textbf{Bit~23 -- padding (wypełniacz):} jest zawsze \texttt{0},
        nie niesie żadnej informacji. Wynika z~rozkładu pól: 1~bit wiodący
        + 7~bitów wykładnika = 8~bitów zajmuje bity 31--24,
        a~mantysa zajmuje bity 22--0 (23~bity).
        Bit~23 zostaje pomiędzy nimi jako martwa strefa.
  \item \textbf{Bity 22--0 -- mantysa $M$:} cyfry za binarnym przecinkiem,
        identyczne jak w~IEEE~754. Dla $9$: $M = 001\underbrace{00\cdots0}_{20}$.
\end{itemize}

Zero kodowane jest jako \texttt{0x00000000}.
Format nie obsługuje liczb ujemnych ani ułamków mniejszych od~1.

Uzasadnienie projektowe:
\begin{itemize}
  \item \textbf{Brak stałej wykładnika}: porównanie dwóch liczb przez
        porównanie słów całkowitoliczbowych; mnożenie i~dodawanie wykładników
        bez korekcji o~bias.
  \item \textbf{Jawny bit wiodący}: brak specjalnego przypadku dla zera
        mantysy; dekoder nie musi niejawnie dodawać~1.
  \item Obydwa uproszczenia skracają czas realizacji operacji arytmetycznych
        w~systemach czasu rzeczywistego.
\end{itemize}

% ===============================================================================
\section*{A.4: Wzajemna konwersja obu formatów}
% ===============================================================================

Ponieważ oba formaty przechowują tę samą znormalizowaną postać $1{,}M \times 2^E$,
a~mantysa (bity~22:0) jest identyczna, jedyną różnicą jest kodowanie wykładnika.

\medskip
\noindent\textbf{Kod wewnętrzny $\to$ IEEE~754:}
\[
  E_b = E + 127, \quad S = 0, \quad M_{\text{IEEE}} = M_{\text{wew}}
\]
\[
  \text{IEEE} = (E_b \ll 23) \mid M
\]

\noindent\textbf{IEEE~754 $\to$ kod wewnętrzny} (tylko dla $S=0$, $E_b \in [127, 254]$):
\[
  E = E_b - 127, \quad \text{bit~31} = 1, \quad M_{\text{wew}} = M_{\text{IEEE}}
\]
\[
  \text{wew} = \texttt{0x80000000} \mid (E \ll 24) \mid M
\]

Warunki konieczne dla konwersji IEEE~754 $\to$ kod wewnętrzny:
\begin{itemize}
  \item $S = 0$ (liczba nieujemna),
  \item $E_b \notin \{0, 255\}$ (liczba znormalizowana, nie Inf/NaN),
  \item $E_b \geq 127$ (wartość $\geq 1{,}0$, bo $E = E_b - 127 \geq 0$),
  \item $E_b \leq 254$ (wykładnik $E \leq 127$, co wynika z~7-bitowego pola).
\end{itemize}

% ===============================================================================
\section*{A.5: Kod binarny naturalny a reprezentacja liczb naturalnych}
% ===============================================================================

Kod binarny naturalny (ang.\ \textit{natural binary code}) to standardowy
zapis pozycyjny przy podstawie~2: liczba $N = \sum_{k=0}^{n-1} b_k \cdot 2^k$.
W~odróżnieniu od kodu Graya czy BCD, kolejne liczby różnią się dowolną liczbą bitów.

Każda liczba naturalna $N \geq 1$ ma jednoznaczną znormalizowaną postać
zmiennoprzecinkową:
\[
  N = 1{,}f \times 2^E, \qquad E = \lfloor \log_2 N \rfloor
\]
Mantysa ułamkowa $M = (N - 2^E) \cdot 2^{23-E}$ (przy $E \leq 23$).
Dla $E > 23$ część ułamkowa jest ucinana: reprezentacja traci precyzję
(dokładnie jak w~IEEE~754 single).

\end{document}
